% !TEX root = ../BeamerTemplate.tex
%%%%%%%%%%%%%%%%%%%%%%%%%%%%%%%%%%%%%%%%%%%%%%%%%%%%%%%%%%%%%%%%%%%%%%%%%%%%%%%%%%

\section{Chapter 1--2 Synthesis}

\begin{frame}[t]{What Chapters 1 and 2 Establish}
\footnotesize
\begin{columns}[T,onlytextwidth]
  \column{0.315\textwidth}
  \begin{block}{1. Why}
    \textbf{38~GHz offers bandwidth, yet a LoS-dominant MIMO channel can remain
    highly correlated.}

    \vspace{0.10cm}
    Can a physical lens improve $\mathbf H$, beyond increasing received power?
  \end{block}

  \column{0.315\textwidth}
  \begin{block}{2. How}
    \textbf{A $20\times20$ transmitarray creates angle-dependent focal responses.}

    \vspace{0.10cm}
    One $3\times3$ channel at 37--39~GHz is evaluated through gain, SVD,
    correlation, conditioning, rank, and capacity.
  \end{block}

  \column{0.315\textwidth}
  \begin{block}{3. What the evidence says}
    \textbf{All three solvers predict directional reinforcement and lower
    correlation.}

    \vspace{0.10cm}
    HFSS predicts stronger relative mode balance; Sionna-RT attributes more of
    the benefit to absolute link gain.
  \end{block}
\end{columns}

\vspace{0.05cm}
\begin{beamercolorbox}[rounded=true,sep=0.60ex,wd=\linewidth]{takeaway}
  \centering
  \textbf{Most defensible conclusion:} in the evaluated $3\times3$ channel, the
  lens reinforces matched links, suppresses most cross-angular coupling, lowers
  correlation, and increases raw capacity.
\end{beamercolorbox}

\vspace{0cm}
\begin{beamercolorbox}[rounded=true,sep=0.45ex,wd=\linewidth]{caution}
  \scriptsize Exact gain, phase, and baseline conditioning remain solver-specific;
  cross-solver agreement is numerical validation, not prototype measurement.
\end{beamercolorbox}
\end{frame}
